\subsection{p-GROUPS}

In this number and the following, the letter \(p\) denotes a prime number (§ 4, no 10, Proposition 16).

DEFINITION 9. A finite group whose order is a power of p is called a p-group.

Let G be a p-group of order \(p^{r}\) . Every divisor of \(p^{r}\) is a power of p ( \(\S 4\) , no. 10, Corollary to Theorem 7). Therefore every subgroup and every quotient group of G is a p-group ( \(\S 4\) , no. 4, Corollary to Proposition 4); the cardinal of every homogeneous space of G is a power of p ( \(\S 5\) , no. 5, Theorem 1).

An extension of a p-group by a p-group is a p-group.

*Examples. (1) A commutative \(p\) -group is isomorphic to a product of cyclic groups \(\mathbf{Z} / p^n\mathbf{Z}\) (cf.~Exercise 19 and also VII, § 4, no. 7, Proposition 7).

\begin{enumerate}
\def\labelenumi{(\arabic{enumi})}
\setcounter{enumi}{1}
\item
  Let \(k\) be a finite field of characteristic \(p\) . The strict triangular group \(T_1(n, k)\) is a \(p\) -group.
\item
  The quaternionic group \(\{\pm 1, \pm i, \pm j, \pm k\}\) is a 2-group (cf.~Exercise 4).*
\end{enumerate}

PROPOSITION 11. Let E be a finite set and G a p-group operating on E. Let \(E^{G}\) denote the set of \(x \in E\) such that gx = x for all \(g \in G\) (the fixed points). Then

\[
\operatorname{Card} \left(\mathrm{E} ^ {\mathrm{G}}\right) \equiv \operatorname{Card} (\mathrm{E}) \quad (\text { mod. } p).
\]

\(E - E^{G}\) is a disjoint union of orbits not reduced to a point. The cardinal of such an orbit is a power of p distinct from \(p^{0} = 1\) and hence a multiple of p.

COROLLARY. Let G be a p-group. If G is not reduced to e, its centre is not reduced to e.

Let G operate on itself by inner automorphisms. The set of fixed points is the centre Z of G. By Proposition 11,

\[
\operatorname{Card} (Z) \equiv \operatorname{Card} (G) \equiv 0 \pmod {p},
\]

whence \(\operatorname{Card}(\mathbf{Z}) \neq 1\) and \(\mathbf{Z} \neq \{e\}\) .

THEOREM 1. Let G be a p-group and \(p^{r}\) its order. There exists a sequence of subgroups of G

\[
\mathbf {G} = \mathbf {G} ^ {1} \supset \mathbf {G} ^ {2} \supset \dots \supset \mathbf {G} ^ {r + 1} = \{e \}
\]

such that \((\mathbf{G},\mathbf{G}^k)\subset \mathbf{G}^{k + 1}\) , \(1\leqslant k\leqslant r\) , and \(\mathbf{G}^k /\mathbf{G}^{k + 1}\) , \(1\leqslant k\leqslant r\) , is cyclic of order \(p\) .

The theorem is true for \(G = \{e\}\) . We prove it by induction on \(\operatorname{Card}(G)\) . Let Z be the centre of G, \(x \neq e\) an element of Z (Corollary to Proposition 11) and \(p^{s}, s \neq 0\) , the order of x. Then \(x^{p^{s-1}}\) is an element of order p and therefore Z contains a subgroup \(G^{r}\) which is cyclic of order p.~By the induction hypothesis, the group \(G' = G/G^{r}\) has a series of subgroups \((\mathbf{G}'^{k})_{1 \leqslant k \leqslant r}\) with the required properties. Let \(\pi: G \to G'\) be the canonical homomorphism. The sequence of subgroups of G defined by \(G^{k} = \pi^{-1}(\mathbf{G}'^{k})\) , \(1 \leqslant k \leqslant r\) , \(G^{r+1} = \{e\}\) is a solution for \(G^{k}/G^{k+1}\) is isomorphic to \(G'^{k}/G'^{k+1}\) for \(1 \leqslant k \leqslant r\) (§4, no. 7, Theorem 4).

COROLLARY. Every p-group is nilpotent.

This follows from no. 3, Proposition 7.

PROPOSITION 12. Let G be a p-group and H a subgroup of G distinct from G. Then:

\begin{enumerate}
\def\labelenumi{(\alph{enumi})}
\item
  The normalizer \(\mathbf{N}_{\mathbf{G}}(\mathbf{H})\) of \(\mathbf{H}\) in \(\mathbf{G}\) is distinct from \(\mathbf{G}\) .
\item
  There exists a normal subgroup N of G of index p in G, which contains H.
\end{enumerate}

Assertion (a) follows from no. 3, Corollary 1 to Proposition 8. We prove (b). By no. 3, Corollary 2 to Proposition 8, there exists a normal subgroup \(\mathbf{N}'\) of \(\mathbf{G}\) containing \(\mathbf{H}\) , distinct from \(\mathbf{G}\) and such that \(\mathbf{G} / \mathbf{N}'\) is commutative. Let \(\mathbf{N}\) be a maximal subgroup distinct from \(\mathbf{G}\) containing \(\mathbf{N}'\) . Then \(\mathbf{N}\) is normal (no. 2, Corollary 3 to Proposition 6) and \(\mathbf{G} / \mathbf{N}\) is a simple commutative \(p\) -group and hence cyclic of order \(p\) ( \(\S 4\) , no. 10, Corollary to Proposition 20).

COROLLARY. Let G be a p-group. Every subgroup of G of index p is normal.

